\section{Robots with ROS2 drivers}
\label{sec-ros2}
\subsection{An opinionated introduction to ROS}
The text in this subsection assumes ROS2, unless stated otherwise.
Furthermore, its purpose is not to provide documentation to ROS,
which is now decent and can be found \href{https://docs.ros.org/}{here}.
Instead, it is to give you, the reader, the minimal necessary context
necessary to decide whether you want to or have to delve into ROS,
and to provide rationale for how SMC is integrated with ROS.

\paragraph{What is ROS}
ROS, which stands for Robotic Operating System, is the biggest
set of open source robotics software.
Its value proposition is to provide a complete software ecosystem
for a robotics platform. The purpose of having a shared ecosystem
is to provide an as uniform as possible API to the various
components of a robotics system --- thus you need to learn 
only one way to program a robot,
instead of multiple frameworks for each component of a robotic system.
This also makes the integration of all components simpler.
Its key component, the feature
that ties everything together, is a multiprocessing and IPC 
framework called RCL. RCL's value proposition is the alleviation
of the need to write explicit multiprocessing and IPC code,
which can be very demanding to implement properly,
and requires a lot of specific knowledge.


Along with RCL, ROS comes with libraries for various components
of a robotic system: path planning (MoveIt), control (ROSControl), navigation (Nav2), 
all kinds of sensor processing, visualization and so on.
Additionally, there are dedicated wrappers around common software building
systems and Linux workflows. For example, ROS packages are build with Colcon
which is a layer on top of CMake.

For more information, the reader is referred to \href{https://docs.ros.org/}{ROS documentation}.


\paragraph{A good situation to use ROS in}
The bigger and more complex the robotic platform in question, the more 
value can be gained by using ROS. This is because if the robotic system
consist of dozens of sensors, multiple computers and possibly hundreds of threads,
management of all this resources becomes practically impossible
if they do not all operate under the same software architecture. 
This is most obvious when working with mobile platform which need
to perform mapping, localization, planning and control,
and make use of ultrasonic distance sensors, Lidars and cameras 
to achieve their tasks.

\paragraph{A situation where ROS is unnecessary}
However, if one works with a simpler robotic system, like a single
manipulator, a force torque sensor and a camera, writing software wrappers
around hardware drivers to conform to a very generic API yields little reward
as these resources can be easily managed individually.
Furthermore, engineers typically work with a few robotic systems
which do not change too often (they are very expensive)
and which they know inside and out. In such cases
adding additional software overhead and dependencies complicates
both software development and maintenance
for very little added functionality.

\paragraph{What is SMC in the context of ROS?}
Within the ROS ecosystem, SMC could be considered as an alternative to ROSControl.
ROSControl of course has better integration with the rest of ROS,
but it has a cumbersome API which makes implementation of controllers time consuming.
On the other hand, since SMC has all workings of an ecosystem,
it could also be considered as an alternative to ROS which targets 
smaller simpler robotic platforms.
It practise, both options work. Which one to use depends 
on the robotic platform in question.

When prototyping whole-body controllers on mobile robots with manipulators,
having ROS run the Nav2 stack to perform mapping and localization,
while using SMC in place of ROSControl gets the best of both worlds.
Upon the completion of the control engineering task, the final controller
can be re-rewritten into ROSControl API to make it more accessible (ROS is popular).

When working with a single manipulator with a few sensors, SMC covers all bases,
and leads to a more flexible workflow due to its leanness.

In the end, one should weight their options and go for what seems like the least
amount of work, which is of course heavily dependent
on one's prior knowledge an experience. As a rule of thumb,
ROS caters more to people with computer science backgrounds,
while SMC caters more to control engineers.

\subsection{SMC-ROS integration approach}
To interact with ROS2 topics programmatically, publishers and subscribers need to be created in a 
ROS \fbox{\texttt{Node}}. This falls under what a RobotManagerXYZ should do.
Getting data from subscribers is only possible through a callback associated with the subscriber object.
Thus, it is also better to publish commands via a dedicated function and timer
to deal with all topics uniformly.

This stands in opposition to AbstractRobotManager's and ControlLoopManager's expected way of working,
but is inelegantly made to work with some ducktaping. 
The main obstacle is that a RobotManagerXYZNode has simultaneously provide functionality of both
ControlLoopManager and RobotManager.
It has to be a ControlLoopManager because it has to publish likely heterogeneous velocity commands via their publishers,
which is meaningless without computing those commands first.
Since only code execution after spin() can happen in topic callbacks,
control needs to be computed in a callback in the Node. There can be only 1 node
as there can be only 1 RobotManager (which is an input to controllerLoops of course).
It also has to read in the current state of the robot via subscription to various topics,
which is the ownership of RobotManager.

This coupling of control signal computation and communication with the robot is exactly
what SMC is trying to avoid --- SMC assumes a fixed robot interface, and swappable controllers.
Controller swapping in ROS is handled via the ROSControl module. The issue with it is its 
verbosity, which slows down quick development and prototyping of controllers (the key
design objective of SMC).
Fortunately, the ducktaping between the two is neither complicated nor too annoying to use,
and the workflow between the two is interoperable where it needs to be.
More on how a RobotManager is made to be a RobotManagerXYZNode can is 
lower in the text. First we introduce main file structure to start a ROS Node.

\subsubsection{Main file structure}
Any file starting a ROS node HAS TO end with the following lines of code:
\begin{minted}[mathescape, linenos]{python}
def main():
    rclpy.init(args=args) # different than smc's args, more on this later
    robot = RobotManagerXYZNode(args_smc, modes_and_loops)
	modes_and_loops = []]

	...

	robot.setModesAndLoops(modes_and_loops)
    executor = MultiThreadedExecutor() # required
    executor.add_node(robot) # required
    executor.spin() # required

if __name__ == "__main__":
    main()
\end{minted}
\paragraph{Hack 1}
The Node can not be used without the executor, and topics can not be interacted with without a Node.
In other words, the robot can only be interacted with from callbacks in the Node.
This in turn means that all control modes and controllers need to be provided to the Node \textbf{before}
executor.spin() is called.
This is done with \fbox{\texttt{modes\_and\_loops}} which is a list of tuples
of control modes and ControlLoopManagers. 

The list is populated like this:
\begin{minted}[mathescape, linenos]{python}
mode_1 = AbstractRobotManager.control_mode.mode_1
loop_manager_controllerX = controllerX(.., args_smc, robot, run=False)
modes_and_loops.append((mode_1, loop_manager_controllerX))
\end{minted}
Hence, instead of immediately running the controller, the output is
a ControlLoopManager instance with the controller and all it's settings.
How the list of controller is used in the RobotManagerXYZNode is described in the following subsection describing
a Node connecting a particular robot in SMC with the ROS2 ecosystem.

\paragraph{Example main file}
An example main file can be found in 
\newline
\fbox{\texttt{examples/ros/mir\_test.py}}.

\subsubsection{Skeleton of ROS2 node for SMC}
Python supports multiple inheritance. This means
we can write a Node which in addition to inheriting from 
\fbox{\texttt{rclpy.node.Node}},
also inherits from both \fbox{\texttt{AbstractRobotXYZRobotManager}}
and \fbox{\texttt{AbstractRealRobotManager}}.
This way we can write a single Node to interact with the robot,
while preserving as much as possible of the SMC architecture.

\paragraph{Hack 2}
Unfortunately one can not use \fbox{\texttt{\_getQ()}}
and \fbox{\texttt{\_getV()}} commands, as SMC does not expects
those to have inputs, while ROS callbacks have to receive and unpack messages.
Hence, the first hack is to leave those as empty, and have the topic callbacks
write to \fbox{\texttt{\_q}} and \fbox{\texttt{\_v}} instead.
At least they remind you that you need to find the appropriate topics
and write their respective callbacks.
(There are other functional design choices, but I felt this was to
be the least convoluted or ugly).

\paragraph{Hack 3}
The second issue is that \fbox{\texttt{sendVelocityCommandToReal}}
does not make sense for the same reason - commands are to be sent via callbacks,
as now the timing for sending messages has to be done via timers in ROS Node's.
The least API-breaking solution is to have a variable
\fbox{\texttt{\_v\_cmd}} which stores the computed velocity command.
\textbf{BIG TODO:} validate that there are no race conditions
on this variable.


\paragraph{RobotManagerNode example}
Look at 
\fbox{\texttt{python/smc/robots/implementations/heron.py}}.


\subsection{Known issues}
\begin{itemize}
		\item Node overrides SIGINT handling from ControlLoopManager, which is in charge 
				of both stopping the robot and saving logs. 
				\textbf{This is a critical issue which needs a proper solution}.
				The current workaround is to rely on the robot's emergency stop button to stop it,
				and to somehow manually trigger log saving (at the moment there is no fully satisfactory solution).
		\item To read the current robot's pose before running a simulation meant to be as close to reality as possible,
				a Node has to be spun (there is no other way to read from a topic programmatically). 
				There is no convenient way to then change topics to those of a ROS2-supported simulator (ex. Gazebo)
				to emulate the topic calls, or to which to a SimulatedRobotManager from within the Node,
				simply running the controllers without publishing. 
				\textbf{This is a critical issue which needs a proper solution. Convenient
				pre-experiment validating simulations are required to safely perform experiments}.
				It is simply too tempting to skip just precautions in the throws of debugging,
				they have to be quick and painless to be actually used.
		\item Node disables opening of windows, blocking plotter from opening a window to plot in.
				Since this is not critical, no solution has been implemented yet.
		\item ROS2 scripts override argument parsing, making it impossible to provide ``normal'' arguments
				to a Python script (which is what's used in SMC). 
				Due to the mutability of the \fbox{\texttt{args}} object ``arguments''
				can be filled in programmatically, but this is an inconvenient hassle.
\end{itemize}

